% ---------------------------------------------------------------------------
% Results (first draft, 2026-10-13).
% Numbers come from figures/results/blind_learned_results.csv (made by
% src/summarize_results.py), the paired tests on figures/scatter{5,7}/scores_mean.csv,
% and the repair counts computed from the saved predictions of the three runs.
% Figures are in figures/paper/: copy them to the figures/ folder of the paper.
% \cite keys match paper/references.bib.
% ---------------------------------------------------------------------------

\section{Results}

\subsection{Measuring the denoising error}

To compare a denoised map with the clean map we need the angle between two
orientations. For a cubic crystal the same orientation can be written in 24
equivalent ways, so a plain difference between two sets of Euler angles is not
a valid measure~\cite{atindama2023restoration,ding2020indexing}. We therefore
use the disorientation angle. Let $\mathbf{q}$ and $\mathbf{q}'$ be the unit
quaternions of two orientations. Their misorientation is
\begin{equation}
  \Delta\mathbf{q} = \mathbf{q}'\,\mathbf{q}^{-1},
  \label{eq:misorientation}
\end{equation}
and the disorientation angle is the smallest rotation angle among the
symmetrically equivalent forms of $\Delta\mathbf{q}$,
\begin{equation}
  \delta(\mathbf{q}, \mathbf{q}') = \min_{\mathbf{s} \in \mathcal{S}}
  \; 2 \arccos \left| \left( \mathbf{s}\, \Delta\mathbf{q} \right)_0 \right|,
  \label{eq:disorientation}
\end{equation}
where $\mathcal{S}$ is the set of the 24 proper rotations of the cubic point
group and $(\cdot)_0$ is the scalar part of a quaternion.

For a map $m$ with the set of pixels $\Omega$ ($|\Omega| = 128 \times 128 =
16{,}384$), denoised map $u_m$ and true map $u^0_m$, the denoising error is the
mean per-pixel disorientation,
\begin{equation}
  e_m = \frac{1}{|\Omega|} \sum_{\mathbf{x} \in \Omega}
        \delta\big(u_m(\mathbf{x}), u^0_m(\mathbf{x})\big).
  \label{eq:error}
\end{equation}
We also compute the same mean over two subsets of the pixels: the
grain-boundary pixels $B_m$, which have at least one of their four neighbours
in a different grain of the true map, and the misindexed pixels $M_m$, which
the noise generator replaced. We call the three measures ALL ($e_m$), BOUNDARY
($e^B_m$) and CORRUPTED ($e^M_m$). Last, the percentage of pixels that are
still grossly wrong is
\begin{equation}
  g_m = \frac{100}{|\Omega|} \sum_{\mathbf{x} \in \Omega}
        \mathbb{1}\Big[\delta\big(u_m(\mathbf{x}), u^0_m(\mathbf{x})\big) > 10^\circ\Big],
  \label{eq:gross}
\end{equation}
where $\mathbb{1}[\cdot]$ is 1 when the condition holds and 0 otherwise.

Each test set has $N = 500$ maps. For each set we report the mean over the
maps, $\bar{e} = \frac{1}{N}\sum_{m=1}^{N} e_m$. Since every map has the same
number of pixels, $\bar{e}$ is also the mean disorientation over all
$N|\Omega|$ pixels of the set. We trained the network three times, with
different random seeds. Unless stated otherwise, the value given for our
method is the mean of the three runs, and the $\pm$ value is the standard
deviation across the runs, which shows how much a result depends on the
random start of training. To compare noise levels we also use the fraction of
the input error that the denoising removed,
\begin{equation}
  R = 1 - \frac{\bar{e}_{\text{out}}}{\bar{e}_{\text{in}}}.
  \label{eq:removed}
\end{equation}

To test whether two methods differ, we use the paired $t$-test on the
per-map errors~\cite{neter1996applied}, as Atindama et
al.~\cite{atindama2023restoration} did. For the same $n$ maps let
$d_m = e^{(1)}_m - e^{(2)}_m$ be the difference between the errors of the two
methods on map $m$, $\bar{d}$ the mean and $S_d$ the standard deviation of the
$d_m$. The test statistic
\begin{equation}
  T = \frac{\bar{d}}{S_d / \sqrt{n}}
  \label{eq:ttest}
\end{equation}
has a $t$-distribution with $n-1$ degrees of freedom. We take $p < 0.05$ as
statistically significant.

\subsection{Training}

Figure~\ref{fig:training_loss} shows the training and validation loss of the
three runs. Each run had 60 epochs of 2{,}437 batches of 16 maps, about
5.8~min per epoch and 5.8~h per run on one NVIDIA H100 GPU. The training and
validation losses fall together and stay close at every epoch, so the network
does not overfit the training maps. Both flatten after about epoch 45. For
each run we kept the checkpoint with the best validation score, which came from
epochs 58, 50 and 58.
% NOTE (for the authors): these three runs chose the checkpoint with a score
% built from MEDIAN errors (the code switched to means on 2026-10-13). Either
% say so in the Methods or retrain before submission.
On the validation maps the gate, the part of the network that marks a pixel as
misindexed, found 98.8\% of the misindexed pixels (recall), and 99.9\% of the
pixels it marked were truly misindexed (precision). The three runs gave the
same values to the third decimal.

\begin{figure}[t]
  \centering
  \includegraphics[width=0.48\textwidth]{figures/training_loss.pdf}
  \caption{Training loss (dashed) and validation loss (solid) of the three
    training runs. The two losses stay close throughout, which shows that the
    network does not overfit the training maps.}
  \label{fig:training_loss}
\end{figure}

\subsection{Maps with mixed noise levels}

We first report one overall result. In the mixed test set, each of the 500
test maps has its own noise level, drawn at random from the same ranges as in
training: $\theta_{\max}$ between $0^\circ$ and $15^\circ$ and $p$ between 0
and 10\%. Table~\ref{tab:mixed} shows the four error measures. The mean
per-pixel error drops from $5.69^\circ$ to $0.589 \pm 0.015^\circ$, so 89.6\%
of the error is removed. The misindexed pixels, which are $40.7^\circ$ from
their true orientation on average, end $2.47^\circ$ from it, and the
percentage of pixels more than $10^\circ$ wrong falls from 11.27\% to 0.23\%.

\begin{table}[t]
  \centering
  \caption{Mean per-pixel errors on the mixed test set (500 maps, each with
    its own random noise level). For our method, mean $\pm$ standard deviation
    over three training runs. $R$ is the fraction of the error removed,
    Eq.~(\ref{eq:removed}).}
  \label{tab:mixed}
  \small
  \begin{tabular}{@{}lccc@{}}
    \toprule
    Measure & Noisy input & Ours & $R$ \\
    \midrule
    ALL ($^\circ$)          & 5.692  & $0.589 \pm 0.015$ & 89.6\% \\
    BOUNDARY ($^\circ$)     & 8.554  & $1.197 \pm 0.020$ & 86.0\% \\
    CORRUPTED ($^\circ$)    & 40.709 & $2.466 \pm 0.032$ & 93.9\% \\
    $>10^\circ$ (\% pixels) & 11.27  & $0.226 \pm 0.002$ & 98.0\% \\
    \bottomrule
  \end{tabular}
\end{table}

\subsection{Effect of the noise level}

Tables~\ref{tab:grid_all} and~\ref{tab:grid_gross} give the results for the
28 test conditions, and Figure~\ref{fig:vs_scatter} plots the errors against
the maximum scatter angle. The standard deviation across the three runs is at
most $0.032^\circ$ for the mean error and 0.006 percentage points for the
percentage of gross errors, far smaller than any of the differences discussed
below.

The error left after denoising grows in proportion to the scatter. With scatter
alone ($p = 0$), the mean error goes from $0.068^\circ$ at
$\theta_{\max} = 1^\circ$ to $0.922^\circ$ at $\theta_{\max} = 15^\circ$, but
the fraction removed stays between 86.5\% and 88.1\% at every level. This is
what we expect from averaging. The network replaces each pixel by a weighted
average of its neighbours in the same grain, and averaging random errors
shrinks them by a fixed factor, whatever their size. So the remaining error is
a fixed fraction of the input error. When misindexing is added, the fraction
removed is larger, between 88.1\% and 93.4\%, because the misindexed pixels
carry a large part of the input error and most of them are repaired.

With misindexing, the percentage of grossly wrong pixels
(Table~\ref{tab:grid_gross}) falls from about $p$ in the input to between
0.10\% and 0.58\% after denoising. At
$\theta_{\max} = 12^\circ$ and $15^\circ$, scatter alone puts 16.7\% and 33.4\%
of the pixels beyond $10^\circ$; after denoising only 0.02\% and 0.05\% remain.

\begin{table*}[t]
  \centering
  \caption{Mean per-pixel error, ALL ($^\circ$), for the 28 test conditions,
    500 maps each: maximum scatter angle $\theta_{\max}$ and percentage of
    misindexed pixels $p$. Ours is the mean of three training runs (standard
    deviation across runs $\leq 0.032^\circ$).}
  \label{tab:grid_all}
  \begin{tabular}{@{}c cc cc cc cc@{}}
    \toprule
    & \multicolumn{2}{c}{$p = 0\%$} & \multicolumn{2}{c}{$p = 3\%$}
    & \multicolumn{2}{c}{$p = 6\%$} & \multicolumn{2}{c}{$p = 9\%$} \\
    \cmidrule(lr){2-3}\cmidrule(lr){4-5}\cmidrule(lr){6-7}\cmidrule(l){8-9}
    $\theta_{\max}$ ($^\circ$) & Input & Ours & Input & Ours & Input & Ours & Input & Ours \\
    \midrule
     1 & 0.500 & 0.068 & 1.708 & 0.128 & 2.913 & 0.196 & 4.121 & 0.273 \\
     3 & 1.500 & 0.186 & 2.679 & 0.248 & 3.853 & 0.317 & 5.031 & 0.396 \\
     5 & 2.501 & 0.301 & 3.649 & 0.365 & 4.793 & 0.437 & 5.942 & 0.518 \\
     7 & 3.501 & 0.418 & 4.619 & 0.485 & 5.733 & 0.560 & 6.852 & 0.645 \\
     9 & 4.501 & 0.538 & 5.589 & 0.609 & 6.673 & 0.688 & 7.762 & 0.777 \\
    12 & 6.002 & 0.724 & 7.044 & 0.803 & 8.084 & 0.891 & 9.127 & 0.989 \\
    15 & 7.502 & 0.922 & 8.500 & 1.013 & 9.494 & 1.111 & 10.492 & 1.222 \\
    \bottomrule
  \end{tabular}
\end{table*}

\begin{table*}[t]
  \centering
  \caption{Percentage of pixels more than $10^\circ$ from their true
    orientation, Eq.~(\ref{eq:gross}), for the same conditions as
    Table~\ref{tab:grid_all}. Ours is the mean of three training runs
    (standard deviation across runs $\leq 0.006$ percentage points).}
  \label{tab:grid_gross}
  \begin{tabular}{@{}c cc cc cc cc@{}}
    \toprule
    & \multicolumn{2}{c}{$p = 0\%$} & \multicolumn{2}{c}{$p = 3\%$}
    & \multicolumn{2}{c}{$p = 6\%$} & \multicolumn{2}{c}{$p = 9\%$} \\
    \cmidrule(lr){2-3}\cmidrule(lr){4-5}\cmidrule(lr){6-7}\cmidrule(l){8-9}
    $\theta_{\max}$ ($^\circ$) & Input & Ours & Input & Ours & Input & Ours & Input & Ours \\
    \midrule
     1 & 0.00 & 0.00 & 2.98 & 0.10 & 5.96 & 0.23 & 8.94 & 0.37 \\
     3 & 0.00 & 0.00 & 2.98 & 0.11 & 5.96 & 0.23 & 8.94 & 0.38 \\
     5 & 0.00 & 0.00 & 2.98 & 0.11 & 5.96 & 0.24 & 8.94 & 0.39 \\
     7 & 0.00 & 0.00 & 2.98 & 0.12 & 5.96 & 0.25 & 8.94 & 0.41 \\
     9 & 0.00 & 0.00 & 2.98 & 0.13 & 5.96 & 0.27 & 8.94 & 0.43 \\
    12 & 16.66 & 0.02 & 19.14 & 0.15 & 21.62 & 0.31 & 24.10 & 0.49 \\
    15 & 33.35 & 0.05 & 35.33 & 0.21 & 37.30 & 0.38 & 39.29 & 0.58 \\
    \bottomrule
  \end{tabular}
\end{table*}

\begin{figure*}[t]
  \centering
  \includegraphics[width=\textwidth]{figures/results_vs_scatter.pdf}
  \caption{Mean per-pixel error against the maximum scatter angle, for each
    percentage of misindexed pixels: (left) all pixels, (middle) grain-boundary
    pixels, (right) misindexed pixels. Dashed lines: noisy input. Solid lines
    with markers: our method, mean $\pm$ standard deviation over three training
    runs (the error bars are smaller than the markers). Logarithmic vertical
    axes.}
  \label{fig:vs_scatter}
\end{figure*}

\subsection{Repair of misindexed pixels}

A misindexed pixel carries no information about its true orientation, so it
must be rebuilt from its neighbours. Over the 21 test conditions with
misindexing, between 91.7\% ($\theta_{\max} = 15^\circ$, $p = 9\%$) and 95.5\%
($\theta_{\max} = 1^\circ$, $p = 3\%$) of the misindexed pixels end within
$5^\circ$ of their true orientation. Their mean error drops from about
$40.7^\circ$ to between $2.02^\circ$ and $3.26^\circ$
(Fig.~\ref{fig:vs_scatter}, right). The repair becomes slightly harder as the
scatter grows. A likely reason is that a misindexed pixel whose random
orientation happens to lie close to the true one is hard to tell apart from a
strongly scattered pixel.

\subsection{Where the remaining errors are}

Figure~\ref{fig:examples} shows one test map under three noise conditions.
The denoised maps are close to the true map of Fig.~\ref{fig:noise_examples}(a),
even at $\theta_{\max} = 15^\circ$ where the grains in the input are hard to
see. In the error maps of the denoised maps, the remaining error follows the
grain boundaries. The numbers confirm this. In every test condition, between
92.5\% and 99.7\% of the pixels that remain more than $10^\circ$ wrong lie on a
grain boundary, although boundary pixels make up only 17\% of a map. With
misindexing, most of these pixels (72--98\%) are misindexed pixels on a
boundary. Such a pixel has neighbours from two grains, and the network
sometimes gives it the orientation of the wrong grain. For the same reason the
BOUNDARY error is always higher than the ALL error
(Fig.~\ref{fig:vs_scatter}, middle).

\begin{figure*}[t]
  \centering
  \includegraphics[width=\textwidth]{figures/denoised_examples.pdf}
  \caption{Test map 1 (Fig.~\ref{fig:noise_examples}(a)) denoised by our
    method (run 1) under three noise conditions, one per row. From left:
    noisy input and denoised map (IPF-Z colours), and the disorientation of
    each to the true map (logarithmic colour scale). The mean per-pixel error
    of each map is given above it. The remaining error follows the grain
    boundaries.}
  \label{fig:examples}
\end{figure*}

\subsection{Comparison with MTEX filters (preliminary)}

So far the MTEX filters~\cite{hielscher2019denoising} have been run on only the
first 10 test maps of two scatter-only conditions, $\theta_{\max} = 5^\circ$
and $7^\circ$, with their default parameters. Table~\ref{tab:mtex} compares
them with our method on these maps. Our method has a lower mean error than the
mean, median, Kuwahara, half-quadratic and infimal-convolution filters on all
10 maps at both levels, and the paired $t$-tests give $p < 0.001$. The spline
filter is the closest: its mean error is $0.403^\circ$ against our
$0.300^\circ$ at $\theta_{\max} = 5^\circ$, and $0.601^\circ$ against
$0.416^\circ$ at $\theta_{\max} = 7^\circ$. However, the spline error changes
a lot from map to map (standard deviation $0.23^\circ$ and $0.30^\circ$,
against $0.009^\circ$ and $0.012^\circ$ for ours). Our method is lower on 6 and
7 of the 10 maps, and with only 10 maps the difference is not statistically
significant ($p = 0.17$ and $p = 0.08$). A full comparison on all 500 test
maps, with filter parameters tuned for each noise level and with misindexed
pixels, is in progress.

\begin{table}[t]
  \centering
  \caption{Mean per-pixel error ALL ($^\circ$) on test maps 1--10, scatter
    only: mean (standard deviation across the 10 maps). Ours is averaged over
    three training runs. $p$-values of the paired $t$-test against our method:
    *** $p < 0.001$; n.s.\ not significant ($p \geq 0.05$).}
  \label{tab:mtex}
  \small
  \begin{tabular}{@{}lll@{}}
    \toprule
    Method & $\theta_{\max} = 5^\circ$ & $\theta_{\max} = 7^\circ$ \\
    \midrule
    Noisy input         & 2.502              & 3.503 \\
    Mean                & 0.986 (0.016) ***  & 1.383 (0.017) *** \\
    Median              & 0.408 (0.014) ***  & 0.575 (0.019) *** \\
    Kuwahara            & 0.977 (0.013) ***  & 1.371 (0.017) *** \\
    Spline              & 0.403 (0.226) n.s. & 0.601 (0.301) n.s. \\
    Half-quadratic      & 1.824 (0.022) ***  & 2.569 (0.030) *** \\
    Infimal convolution & 1.916 (0.015) ***  & 2.903 (0.022) *** \\
    \textbf{Ours}       & \textbf{0.300} (0.009) & \textbf{0.416} (0.012) \\
    \bottomrule
  \end{tabular}
\end{table}
